\clearpage
\fancyhead[L]{\small GAUSS TO BODE.}
\SourceLeaf{454}

\begin{center}
{\large GAUSS TO BODE.}
\end{center}

\begin{center}
\rule{6cm}{0.4pt}\\
Astronomisches Jahrbuch für 1827. pp. 128..130. Berlin 1824.\\
\rule{6cm}{0.4pt}
\end{center}

\begin{flushright}
Göttingen, 21 April 1824.
\end{flushright}

In all these observations the comet was so faint that it could not tolerate any illumination of the wires. The comet and the reticle could therefore be made visible only by alternately blocking and admitting the light. In these circumstances the observations, especially those of right ascension, could be made with only limited precision. The times given refer to the mean of the observed, or rather estimated, wire crossings. Thus on days when observations were not made at every wire they differ somewhat from the true time of culmination.

\begin{center}
\begingroup\small\setlength{\tabcolsep}{4pt}
\begin{tabular}{lr|c|c|c}
\multicolumn{2}{c|}{1824} & \shortstack{Mean time\\in Göttingen} & \shortstack{Right\\ascension} & Declination \\
\hline
Feb. & 31 & $13^{\mathrm u}\,28^{\mathrm m}\,20^{\mathrm s}$ & $152^\circ\,16'\,29''$ & $72^\circ\,38'\,35''$ N. \\
Jan. & 3 & $11\quad55\quad7$ & $131\quad23\quad10$ & $69\quad14\quad9$ \\
 & 7 & $10\quad38\quad3$ & $116\quad12\quad50$ & $63\quad14\quad56$ \\
 & 8 & $10\quad24\quad5$ & $113\quad50\quad49$ & $61\quad46\quad15$ \\
 & 19 & $8\quad52\quad23$ & $101\quad49\quad31$ & $48\quad48\quad51$ \\
 & 20 & $8\quad48$ & ---\quad---\quad--- & $47\quad55$ \quad : \\
 & 27 & $8\quad11\quad46$ & $99\quad26\quad51$ & $42\quad48\quad49$ \\
 & 28 & $8\quad6\quad47$ & $99\quad18\quad8$ & $42\quad12\quad46$ \\
March & 2 & $7\quad53\quad59$ & $99\quad0\quad39$ & $40\quad32\quad3$ \\
\end{tabular}
\endgroup
\end{center}

I have not yet reduced some earlier observations I made away from the meridian with a ring micrometer or heliometer, because I must first determine for myself the comparison stars that are absent from the Histoire Céleste.

Dr. \textsc{Schmidt}, a young man with outstanding mathematical knowledge, calculated the following parabolic elements from these and some other observations.

\begin{center}
\begingroup\small
\begin{tabular}{p{.51\linewidth}p{.41\linewidth}}
Longitude of the node \mbox{.\,.\,.\,.\,.\,.\,.\,.\,.\,.\,.\,.} & $303^\circ\,3'\,51''3$ \\
Inclination of the orbit \mbox{.\,.\,.\,.\,.\,.\,.\,.\,.\,.\,.\,.} & $76\quad12\quad14.0$ \\
Longitude of perihelion \mbox{.\,.\,.\,.\,.\,.\,.\,.\,.\,.\,.} & $274\quad34\quad14.4$ \\
Passage through perihelion \mbox{.\,.\,.\,.} & 1823 December 9.47459 Göttingen time. \\
Logarithm of the least distance. \mbox{.\,.\,.\,.\,.\,.} & 9.3553041 \\
\multicolumn{2}{l}{Retrograde motion.} \\
\end{tabular}
\endgroup
\end{center}

\Needspace{23\baselineskip}
Comparing these elements with the observations reported so far gave the following differences.

\Needspace{12\baselineskip}
\SourceLeaf{455}


\begin{center}
\begin{tabular}{lr|r|r|l}
 & & \shortstack{in right\\ascension} & \shortstack{in\\declination} & Observer \\
\hline
Jan. & 1 & 0 & 0 & Bouvard \\
 & 3 & $+12.5$ & $-22.6$ & Nicolai \\
 & 4 & $+13.6$ & $+1.6$ & Nicolai \\
 & 5 & $+34.5$ & $-6.5$ & Nicolai \\
 & 11 & $+0.2$ & $+9.5$ & Olbers \\
 & 14 & $+27.5$ & $+0.3$ & Olbers \\
 & 24 & $+48.6$ & $+3.4$ & Schumacher \\
 & 25 & $+85.0$ & $-3.3$ & Schumacher \\
 & 26 & $+44.0$ & $+11.0$ & Soldner \\
 & 27 & $+33.0$ & $+10.4$ & Soldner \\
 & 30 & $-26.5$ & $+20.7$ & Schumacher \\
 & 31 & $-36.7$ & $+25.5$ & Soldner \\
 & 31 & $+24.9$ & $+14.5$ & Gauss \\
Feb. & 2 & $+963.4$ & $-5.7$ & Schumacher \\
 & 3 & $-28.9$ & $+12.0$ & Schumacher \\
 & 3 & $-40.6$ & $-8.6$ & Gauss \\
\end{tabular}
\end{center}

\begin{center}
\begin{tabular}{lr|r|r|l}
 & & \shortstack{in right\\ascension} & \shortstack{in\\declination} & Observer \\
\hline
Feb. & 7 & $-4.5$ & $+18.5$ & Gauss \\
 & 8 & $-40.6$ & $-2.0$ & Gauss \\
 & 11 & $-18.4$ & $-22.0$ & Schumacher \\
 & 13 & $-10.4$ & $-1.2$ & Schumacher \\
 & 18 & $-2.6$ & $-66.1$ & Olbers \\
 & 19 & $-1.0$ & $+1.0$ & Gauss \\
 & 19 & $+26.0$ & $-26.7$ & Olbers \\
 & 21 & $+36.2$ & $-38.8$ & Olbers \\
 & 27 & $+27.8$ & $-0.2$ & Gauss \\
 & 27 & $+21.6$ & $-34.5$ & Olbers \\
 & 28 & $+21.4$ & $-22.7$ & Gauss \\
 & 28 & $+0.4$ & $-11.4$ & Olbers \\
March & 2 & $-41.3$ & $+9.8$ & Gauss \\
 & 2 & $+40.1$ & $-34.9$ & Olbers \\
 & 5 & $+7.3$ & $-83.3$ & Olbers \\
 & 19 & $+56.3$ & $+36.4$ & Olbers \\
\end{tabular}
\end{center}


There seems to have been an error of one minute of time in the right ascension of the Altona observation of 2 February.
